\documentclass[14pt,a4paper,oneside]{extbook}
\usepackage[utf8]{vietnam}
\usepackage{amsmath,amssymb,mathrsfs,fancyhdr,enumerate,multirow,makecell}
\usepackage{tkz-euclide,tikz-3dplot,tikz,tkz-tab}
\usetikzlibrary{arrows.meta,calc,intersections,angles,quotes,backgrounds,shapes.geometric,patterns,through,shadows,positioning}
\usepackage{yhmath}
\usepackage[top=0.7cm, bottom=0.7cm, left=1.0cm, right=1.0cm, headheight=16pt, headsep=3mm, footskip=5mm]{geometry}
\usepackage[hidelinks,unicode]{hyperref}
\usepackage{qrcode}
\usepackage[many]{tcolorbox}
\usepackage{fontawesome}

\pagestyle{fancy}
\fancyhf{}
\fancyhead[L]{\footnotesize\bfseries\color{teal!85!black}\faBook~SỔ TAY KIẾN THỨC TOÁN 5}
\fancyhead[R]{\footnotesize\bfseries\color{blue!85!black}Lớp Toán Thầy Lĩnh -- loptoanthaylinh.vn}
\fancyfoot[C]{\footnotesize\bfseries\color{gray!80!black}-- Sổ tay kiến thức: Bài 31. Ôn tập các phép tính với số thập phân --}
\renewcommand{\headrulewidth}{0.4pt}

\begin{document}

% --- BANNER TIÊU ĐỀ BÀI HỌC ---
\noindent
\begin{tcolorbox}[
  enhanced,
  colback=teal!8!white,
  colframe=teal!80!black,
  arc=3mm,
  boxrule=1.1pt,
  width=\linewidth,
  halign=center,
  top=2.2mm, bottom=2.2mm,
  left=3mm, right=3mm,
  drop fuzzy shadow
]
  {\large\bfseries\color{teal!95!black}\faBookmark\ BÀI 31. ÔN TẬP CÁC PHÉP TÍNH VỚI SỐ THẬP PHÂN}\\[2pt]
  {\scriptsize\color{black!75}\itshape Rèn luyện thành thạo Cộng, Trừ, Nhân, Chia số thập phân và tính giá trị biểu thức trên 1 trang A4}
\end{tcolorbox}

\vspace{3.2mm}
% --- KHỐI 1: BẢNG TỔNG KẾT QUY TẮC TÍNH ---
\noindent
\begin{tcolorbox}[
  enhanced,
  title={\footnotesize\bfseries\color{white}\faCalculator\ 1. Điểm cốt tử của từng phép tính số thập phân},
  colback=blue!2!white,
  colframe=blue!75!black,
  arc=2.5mm, boxrule=0.8pt, top=1.8mm, bottom=1.8mm, left=2.5mm, right=2.5mm
]
  \footnotesize
  $\bullet$ \textbf{Phép cộng \& trừ:} \textbf{Dấu phẩy thẳng hàng tuyệt đối}. Thêm chữ số $0$ vào số bị trừ nếu thiếu chữ số thập phân.\\
  $\bullet$ \textbf{Phép nhân:} Đặt tính như nhân số tự nhiên. Đếm \textbf{tổng số chữ số TP} của cả hai thừa số để tách ở tích.\\
  $\bullet$ \textbf{Phép chia:} Chuyển dấu phẩy ở SC sang tận cùng bên phải, dịch phẩy ở SBC bấy nhiêu hàng rồi chia cho số tự nhiên.
\end{tcolorbox}

\vspace{3mm}

% --- KHỐI 2: TÍNH NHANH VÀ TÌM X ---
\noindent
\begin{minipage}[t]{0.52\textwidth}
  \begin{tcolorbox}[
    enhanced,
    title={\footnotesize\bfseries\color{white}\faBolt\ 2. Kỹ thuật ghép số tính nhanh},
    colback=teal!2!white,
    colframe=teal!75!black,
    arc=2.5mm, boxrule=0.8pt, top=1.8mm, bottom=1.8mm, left=2.2mm, right=2.2mm
  ]
    \footnotesize
    $\bullet$ Ghép cặp tạo số tròn: $(2,4 + 7,6) = 10$; $(1,25 \times 8) = 10$; $(2,5 \times 4) = 10$.\\
    $\bullet$ Đặt thừa số chung: $a \times b + a \times c = a \times (b + c)$.\\
    \textit{VD:} $6,4 \times 3,8 + 6,4 \times 6,2 = 6,4 \times (3,8 + 6,2) = 6,4 \times 10 = \mathbf{64}$.
  \end{tcolorbox}
\end{minipage}%
\hfill
\begin{minipage}[t]{0.46\textwidth}
  \begin{tcolorbox}[
    enhanced,
    title={\footnotesize\bfseries\color{white}\faSearch\ 3. Tìm thành phần chưa biết},
    colback=purple!2!white,
    colframe=purple!70!black,
    arc=2.5mm, boxrule=0.8pt, top=1.8mm, bottom=1.8mm, left=2mm, right=2mm
  ]
    \footnotesize
    $\bullet$ $x + a = b \Rightarrow x = b - a$.\\
    $\bullet$ $x - a = b \Rightarrow x = b + a$.\\
    $\bullet$ $a - x = b \Rightarrow x = a - b$.\\
    $\bullet$ $x \times a = b \Rightarrow x = b : a$.\\
    $\bullet$ $x : a = b \Rightarrow x = b \times a$.\\
    $\bullet$ $a : x = b \Rightarrow x = a : b$.
  \end{tcolorbox}
\end{minipage}

\vspace{3mm}

% --- KHỐI 3: SƠ ĐỒ TƯ DUY ---
\noindent
\begin{tcolorbox}[
  enhanced,
  title={\footnotesize\bfseries\color{white}\faSitemap\ Sơ đồ tư duy: Bốn phép tính số thập phân},
  colback=yellow!2!white,
  colframe=teal!85!black,
  arc=3mm, boxrule=0.9pt, top=1.8mm, bottom=1.8mm, left=1.5mm, right=1.5mm,
  halign=center
]
  \begin{tikzpicture}[
    scale=0.92,
    every node/.style={transform shape},
    root/.style={rectangle, rounded corners=5pt, fill=teal!85!black, text=white, font=\bfseries\footnotesize, inner sep=5pt, align=center},
    b1/.style={rectangle, rounded corners=4pt, fill=blue!12, draw=blue!75!black, text=blue!90!black, font=\bfseries\scriptsize, inner sep=4.5pt, align=center},
    b2/.style={rectangle, rounded corners=4pt, fill=green!15, draw=green!65!black, text=green!80!black, font=\bfseries\scriptsize, inner sep=4.5pt, align=center}
  ]
    \node[root] (core) at (0, 0.75) {\faCalculator\ TÍNH TOÁN SỐ THẬP PHÂN};
    \node[b1] (n1) at (-4.7, -0.45) {\textbf{CỘNG / TRỪ SỐ THẬP PHÂN}\\[2pt]Thẳng hàng dấu phẩy, thêm $0$ nếu thiếu};
    \node[b2] (n2) at (4.7, -0.45) {\textbf{NHÂN / CHIA SỐ THẬP PHÂN}\\[2pt]Quy tắc đếm chữ số TP \& dịch dấu phẩy};

    \draw[teal!85!black, thick, -{Stealth[length=2mm]}] (core) -| (n1);
    \draw[teal!85!black, thick, -{Stealth[length=2mm]}] (core) -| (n2);
  \end{tikzpicture}
\end{tcolorbox}

\vspace{3mm}

% --- KHỐI 4: NHẨM NHANH & BẪY TRẮC NGHIỆM ---
\noindent
\begin{minipage}[t]{0.54\textwidth}
  \begin{tcolorbox}[
    enhanced,
    title={\footnotesize\bfseries\color{white}\faBolt\ Nhẩm nhanh với $10, 100$ và $0,1, 0,01$},
    colback=green!3!white,
    colframe=green!60!black,
    arc=2.5mm, boxrule=0.8pt, top=1.8mm, bottom=1.8mm, left=2.2mm, right=2.2mm
  ]
    \footnotesize
    $\bullet$ Nhân $10, 100\dots$ hoặc Chia $0,1; 0,01\dots$: \textbf{Dịch phẩy sang phải}.\\
    $\bullet$ Chia $10, 100\dots$ hoặc Nhân $0,1; 0,01\dots$: \textbf{Dịch phẩy sang trái}.
  \end{tcolorbox}
\end{minipage}%
\hfill
\begin{minipage}[t]{0.44\textwidth}
  \begin{tcolorbox}[
    enhanced,
    title={\footnotesize\bfseries\color{white}\faExclamationTriangle\ Mẹo tránh bẫy trắc nghiệm},
    colback=red!2!white,
    colframe=red!70!black,
    arc=2.5mm, boxrule=0.8pt, top=1.8mm, bottom=1.8mm, left=2.2mm, right=2.2mm
  ]
    \footnotesize
    $\bullet$ Trong phép chia có dư: Phải dóng thẳng cột theo dấu phẩy của số bị chia ban đầu để kết luận số dư.\\
    $\bullet$ Thứ tự thực hiện tính ngoặc trước, nhân chia trước.
  \end{tcolorbox}
\end{minipage}

\end{document}
